\section{Limitations, Threats to Validity and Safety}
\label{sec:limits}

\subsection{Limitations}
\emph{Geography.} The maps carry no georeferencing. The reference point $(\varphi_0,\lambda_0)$ is a command-line value whose default is Varanasi city for every map, and $\rho=2$~m/px is a configuration constant that is not derived from the image. Section~\ref{sec:geovalid} shows that both are wrong by one to two orders of magnitude for the supplied maps. Even with correct values, the conversion assumes north-up, square pixels of uniform scale, whereas web-Mercator basemaps change scale with latitude \cite{battersby2014implications}. The spherical constant $L$ introduces a scale error of up to 0.7\%, and the flat-earth approximation degrades over hundreds of kilometres (7.7~km at the Assam extent). The rounding of the display height adds an anisotropy of up to 0.06\%. Correcting it would require separate horizontal and vertical scales in the conversion, which is a change to the method.

\emph{Perception.} Segmentation extracts annotation colours from map products rather than detecting water in imagery. It depends on operator clicks and colour margins, and its accuracy is unknown because no ground-truth masks exist. On thematic maps with white backgrounds it can fail catastrophically (83\% flood for Assam). The hue margin is inert for red annotations.

\emph{Prediction.} The spread rule is a heuristic proxy, not a hydrological model \cite{teng2017flood}. It ignores terrain, drainage and flow, works in display pixels so that its physical reach depends on image resolution, rounds horizons down to whole hours, and treats Open-Meteo's current precipitation, which is a sum over the preceding reporting interval \cite{openmeteo_docs}, as an hourly rate. Its forecast skill has not been measured.

\emph{Selection and ordering.} Drop points are boundary vertices with no terrain, structure, population or landing-safety assessment. They are chosen relative to a centroid polyline in discovery order rather than to the final route, and they usually lie inside $O$ (98 endpoint snaps in the two dot maps). The ordering heuristic is up to 44.5\% longer than optimal in our random instances and ignores obstacles and vehicle endurance.

\emph{Planning.} The $5\times5$ nearest-neighbour downsampling inspects one pixel per block. Snapping can move endpoints into enclosed pockets, and the resulting no-path legs are flown straight across flooded or predicted areas (14 legs in our runs). D*~Lite is rebuilt per leg, so its incremental capability is unused and it is 61--112 times slower than A* on the same task. Routes depend on display width because the grid cell is defined in display pixels.

\emph{Mission design.} The servo is set to 2000~$\mu$s at every drop and never reset, so a release mechanism that requires a toggle would release only once. LAND after RTL is redundant on ArduCopter \cite{ardupilot_missioncmds}. Waypoint yaw 0 would point a PX4 vehicle north at every waypoint \cite{px4_missions}. Servo channel 9 and the 100/10~m altitudes are assumptions, and altitudes are relative to home without terrain following. Mission length is not checked against the flight controller's mission capacity, which matters for missions with more than 1\,000 items such as Kanpur's. The route length is not checked against vehicle endurance; at the estimated real scale of Varanasi it would be 263~km.

\subsection{Threats to Validity}
\emph{Internal validity.} Correctness evidence is strongest for the geospatial and mission layers, which were verified against an independent oracle and a validator that were themselves tested and mutation-checked. The perception, prediction and planning stages are evaluated behaviourally, not against ground truth. All experiments use a single simulated operator and a single weather setting, so operator and weather variability are only partly explored by the sensitivity study. With calm weather the predicted mask equals $M$, so the ablation's current-only condition also describes that case. Runtimes are single measurements on one machine.

\emph{Construct validity.} Route exposure measures geometric overlap with masks, not risk to the vehicle or to people. ``No-path legs'' and ``validator errors'' measure planner completeness and format conformance, not mission success. The landmark-based scale estimate is a plausibility check whose label positions were read off by eye.

\emph{External validity and dataset validity.} The three maps are convenience samples of one style of news graphic and one thematic map, all from India; they are not a benchmark. Their copyright belongs to third parties, and the actual flood observations behind them are not available. The results cannot be generalised to UAV or satellite imagery, to other annotation colours, or to other regions.

\emph{Literature coverage.} The review was targeted rather than systematic: candidates were identified by keyword searches and verified against DOI registries. Statements that a capability is ``not reported'' are therefore bounded by the reviewed set.

\emph{Geographic validity.} For the supplied maps, the generated coordinates do not correspond to the flooded places (Section~\ref{sec:geovalid}). The conversion is correct for correctly specified inputs; obtaining such inputs is outside the current system.

\emph{Operational validity.} Mission semantics were checked against documentation only. No ground station has loaded the missions, and no SITL, HITL or flight test has been performed, so the behaviour of a real autopilot and vehicle is unverified.

\emph{Reproducibility.} The code, inputs, evaluation scripts and raw results are versioned with the paper, and all experiments are deterministic except runtimes. Python dependency versions are recorded but not pinned in the project's requirements, which may affect future reproduction.

\subsection{Safety and Operational Considerations}
\adapt{} is \emph{not} flight-ready and must not be used to fly over people or property. Any operational use would require, at minimum: georeferenced inputs and an independent check of every coordinate against base maps; a ground-station load and SITL rehearsal of the exact mission; geofences, altitude limits and failsafes for link loss, battery and GNSS degradation configured on the autopilot; a trained pilot able to take control; verification that drop points are clear of people, structures and water hazards; a payload-release mechanism tested on the ground with the configured channel and PWM; and compliance with local aviation regulations. The software design supports some of these practices: it rejects invalid coordinates before writing a file, it never leaves a partial or stale mission at the output path, and it logs every no-path leg so that an operator can inspect straight segments before flight. It cannot, however, detect that a valid-looking coordinate is in the wrong place.
